% JAC1964-20261121. New finite-place deduction; no reliance on an integer-k
% assertion or on reading node thickness from a stable intersection graph.
% /root/programme_audit; GPT-6 Astra Pro; effort setting not exposed.

\subsection*{JAC-P10. The finite-place Hodge lattice of the Fermat quartic}

\begin{proposition}[JAC-P10]
\label{jac:P10}
Let $C:x^4+y^4+z^4=0$ over $\overline{\mathbf Q}$, let
$J=\operatorname{Jac}(C)$, and let $E$ be any elliptic curve with
$j(E)=1728$. For stable Faltings heights, in any common normalization,
\[
 h_F(J)-3h_F(E)\ =\ \frac14\log 2.                 \tag{JAC.P10.1}
\]
More precisely, for the degree-eight isogeny
$\Phi:J\to E_x\times E_y\times E_z$ induced by the three coordinate
sign quotients, the normalized valuation of the determinant of
pullback on the good-reduction Hodge lattices at every place above
$2$ is exactly $7/4$, when the valuation is normalized by $\nu(2)=1$.
\end{proposition}

\begin{proof}
The statement is unchanged by a finite extension of the number field.
We may therefore work locally over a finite extension of $\mathbf Q_2$
on which the curve is stable, its stable components and nodes are split,
and the three elliptic quotients have good reduction. We may make
further finite extensions in the local calculation: normalized
determinant valuations of maps between good-reduction Hodge lattices
are invariant under such extension.

\paragraph{A smooth chart at an elliptic tail.}
Put $u=x/z$, $v=y/z$. The quotient by $u\mapsto-u$ has affine
equation $w^2+v^4+1=0$, with $w=u^2$.
Liu, \emph{Algebraic Geometry and Arithmetic Curves},
Example~10.4.39, printed pp.~548--549, gives constants
$\xi_8,\alpha,\beta,\gamma$ over a finite extension such that
\[
 \begin{gathered}
 \xi_8v=1+\frac1{t+\alpha},\qquad
 w=\frac{2s+\beta t+\gamma}{(t+\alpha)^2},\\
 s^2+(\beta t+\gamma)s=t^3,\qquad
 \nu(\alpha)=-\frac14,\quad
 \nu(\beta)=\frac12,\quad \nu(\gamma)=0.
                                                        \tag{JAC.P10.2}
 \end{gathered}
\]
The last equation is a smooth good-reduction elliptic model: its
reduction is $s^2+\bar\gamma s=t^3$.
Choose $a=\sqrt\gamma/\alpha$, put $b=\alpha^{-2}$, and substitute
$u=a+bU$. Thus $\nu(a)=1/4$ and $\nu(b)=1/2$.
The curve is described, on a suitable affine open, by the two
integral equations
\[
 \begin{split}
 G={}&s^2+(\beta t+\gamma)s-t^3=0,\\
 F={}&\gamma t^2+(2\gamma\alpha-\beta\alpha^2)t
       -2\alpha^2s
       +2a(t+\alpha)^2U
       +\alpha^{-2}(t+\alpha)^2U^2=0.
 \end{split}                                             \tag{JAC.P10.3}
\]
Indeed $F$ is $\alpha^2$ times
$u^2(t+\alpha)^2-(2s+\beta t+\gamma)$, after using
$a^2=\gamma/\alpha^2$. Every displayed coefficient is integral.
Its reduction is
\[
 \bar G=s^2+\bar\gamma s-t^3,
 \qquad
 \bar F=U^2+\bar\gamma t^2-\overline{\beta\alpha^2}t,
                                                        \tag{JAC.P10.4}
\]
where both $\bar\gamma$ and $\overline{\beta\alpha^2}$ are nonzero.
The Jacobian determinant of $(F,G)$ with respect to $(t,s)$ is a
unit. Hence this chart is smooth over the valuation ring and $dU$
generates its relative differentials.

The reduction of $F$ gives a purely inseparable extension of degree
two of the elliptic function field defined by $\bar G$. It is
nontrivial because the derivative of
$\bar\gamma t^2-\overline{\beta\alpha^2}t$ is nonzero.
It therefore has genus one. This smooth chart supplies the generic
point of one elliptic component $\Gamma_x$ in the stable reduction.
To justify persistence, compare a proper completion of this chart
with the stable model through a common regular resolution. A
component that disappears over a nodal stable model belongs to the
rational exceptional chains resolving nodes or to rational curves
introduced by point blowups. The genus-one divisorial valuation
just constructed cannot be among them. Equivalently, this is the
reduced-special-fibre case of Liu, Proposition~10.4.37, p.~547.
Here
\[
 \nu_{\Gamma_x}(u)=\frac14,\qquad
 \nu_{\Gamma_x}(v)=0,
 \qquad
 \nu_{\Gamma_x}\!\left(\frac{du}{v^3}\right)=\frac12.
                                                        \tag{JAC.P10.5}
\]
The last equality follows directly from $du=\alpha^{-2}dU$.
The cyclic permutation $(x,y,z)\mapsto(y,z,x)$ supplies the other
two elliptic tails. Their distinction follows from their different
valuations of $u,v$. Since these three genus-one components already
account for arithmetic genus three, every other component is
rational and the dual graph is a tree. A stable rational vertex has
valence at least three. The order-three cyclic symmetry permutes
the three positive-genus vertices, so none can be the middle vertex
of a path connecting them. The tree has just one rational branching
vertex and the three genus-one leaves; any additional rational
vertex would have valence at most two. This is the stable star,
also stated in Liu's example.

\paragraph{Three integral stable differentials.}
Write
\[
 \eta_0=\frac{du}{v^3},\qquad
 \eta_1=u\eta_0,\qquad \eta_2=v\eta_0,
 \qquad \kappa=(1+u+v)\eta_0.
                                                        \tag{JAC.P10.6}
\]
These are holomorphic differentials on the smooth plane quartic;
$\eta_0,\eta_1,\eta_2$ form a basis. At $\Gamma_x$ their vertical
valuations are respectively $1/2,3/4,1/2$.
Since $\nu(\xi_8+1)=1/4$ and $\nu((t+\alpha)^{-1})=1/4$,
formula (JAC.P10.2) gives $\nu_{\Gamma_x}(1+v)\ge1/4$.
Together with $\nu_{\Gamma_x}(u)=1/4$ this yields
\[
 \nu_{\Gamma_x}(\kappa)\ge\frac34.
\]
The cyclic permutation permutes $\eta_0,\eta_1,\eta_2$ and fixes
$\kappa$. Thus each $\eta_i$ has valuation at least $1/2$ on
every elliptic tail, and $\kappa$ has valuation at least $3/4$
on every elliptic tail.

We spell out why checking the tails suffices for integrality.
Let $M=H^0(\mathcal C,\omega_{\mathcal C/R})$ for the stable model.
This is a free rank-three lattice and its formation commutes with
residue-field base change. For a compact-type stable curve with
this star graph, restriction gives an injection
\[
 H^0(\mathcal C_s,\omega_{\mathcal C_s})
       \hookrightarrow\bigoplus_{i=x,y,z}
                           H^0(\Gamma_i,\omega_{\Gamma_i}).
                                                        \tag{JAC.P10.7}
\]
Indeed a dualizing differential has opposite residues on the two
branches of each node; the single residue on each elliptic tail
is zero, so all three residues on the rational component are zero,
and that component has no holomorphic differentials.
Suppose a generic holomorphic differential has nonnegative vertical
valuation on all three tails but is not in $M$. Multiplying it by
the smallest positive power of a uniformizer that puts it in $M$
would give a primitive section whose reduction vanishes on all
three tails. Equation (JAC.P10.7) and base change imply that its
reduction is zero, contradicting primitivity. This proves the claim.

Choose constants $c_2,c_3$ with valuations $1/2,3/4$, respectively.
The preceding observations prove that
\[
 f_1=\eta_0/c_2,\qquad f_2=\eta_1/c_2,
 \qquad f_3=\kappa/c_3
 \quad\text{all belong to }M.                            \tag{JAC.P10.8}
\]
In fact these three differentials form a basis of $M$. To verify
this stronger assertion, take $c_3=c_2c_1$ with $\nu(c_1)=1/4$.
Putting $t=s=0$ in the elliptic-coordinate identity gives
\[
 \gamma^2=(\alpha+1)^4-\alpha^4
   =1+4\alpha^3+6\alpha^2+4\alpha.
\]
Every nonconstant term has positive valuation; hence
$\bar\gamma=1$, and also $\overline{\sqrt\gamma}=1$.
At the generic point of $\Gamma_x$ we have the exact expansion
\[
 \begin{split}
 1+u+v={}&(1+\xi_8^{-1})
       +\frac{\sqrt\gamma+\xi_8^{-1}}{\alpha}
       +\alpha^{-2}U\\
 &+\xi_8^{-1}\bigl((t+\alpha)^{-1}-\alpha^{-1}\bigr).
 \end{split}                                             \tag{JAC.P10.8a}
\]
The first term has valuation exactly $1/4$. The second has valuation
strictly greater than $1/4$, because both terms in its numerator
reduce to one; the last two have valuation at least $1/2$.
Thus $(1+u+v)/c_1$ reduces to a nonzero constant $c$.
Choose the nonzero reduction of $\eta_0/c_2$ as a basis of
holomorphic differentials on $\Gamma_x$, and transport it cyclically
to the other tails. Since $\bar v=1$ and $\nu(u)>0$
on $\Gamma_x$, the restrictions of $f_1,f_2,f_3$ to the three tails
are, up to the choice of cyclic ordering, the rows of
\[
 \begin{pmatrix}
 1&0&c\\ 0&1&c\\ 1&1&c
 \end{pmatrix}.
                                                        \tag{JAC.P10.8b}
\]
Its determinant is $c\ne0$ in residue characteristic two.
Restriction in (JAC.P10.7) is an isomorphism, since both sides have
dimension three. The reductions of $f_1,f_2,f_3$ are therefore a
basis of $M/\mathfrak m_R M$, and Nakayama's lemma makes them a
basis of $M$. Their change-of-basis determinant from
$(\eta_0,\eta_1,\eta_2)$ has valuation $-7/4$.
Consequently the lattice spanned by the $\eta_i$ has index in $M$
of normalized valuation exactly $7/4$.

\paragraph{Identification of the pullback lattice and the height.}
An invariant differential on the good model of the first elliptic
quotient is
\[
 \frac{dt}{2s+\beta t+\gamma}.
\]
By (JAC.P10.2), its pullback is a unit multiple of $dv/w=-\eta_1$.
The other two quotient pullbacks are unit multiples of
$-\eta_2$ and $-\eta_0$. Thus the lattice just computed is precisely
the image of the Hodge lattice under
$\Phi^*:H^0(E_x\times E_y\times E_z,\Omega^1)\to H^0(J,\Omega^1)$.
The identification with $M$ follows from the relative Picard
description of a compact-type Jacobian and relative Serre duality.

For completeness, $\deg\Phi=8$. Each quotient map has degree two;
its norm followed by pullback is $[2]$, and the mixed norm--pullback
maps are zero because the three spaces have different characters
under the coordinate-sign group. With principal polarizations,
$\Phi\Phi^t=[2]$ on the product, and hence
$(\deg\Phi)^2=2^6$.

Let $r_v$ be the unnormalized length of the cokernel of the Hodge
pullback over a number field $L$ where all varieties have good
reduction, and let $e_v=\operatorname{ord}_v(2)$.
The local calculation gives $r_v=7e_v/4$ for $v\mid2$.
At every other finite place the isogeny has invertible degree,
so its Hodge pullback is an isomorphism.
The isogeny height formula, with its sign written explicitly, is
\[
 h_F(E_x\times E_y\times E_z)-h_F(J)
 =\frac12\log8-
   \frac1{[L:\mathbf Q]}\sum_{v\mid2} r_v\log N(v).
                                                        \tag{JAC.P10.9}
\]
One may obtain the sign directly: pullback multiplies the squared
archimedean norm of a top differential by $\deg\Phi$, while the
finite lattice inclusion has the displayed cokernel lengths.
It also agrees with R\'emond, \emph{Propri\'et\'es de la hauteur
de Faltings}, Proposition~2.1 and its proof, author PDF pp.~2--3,
and the differential formulation in \S3, p.~4.
As $\sum_{v\mid2}e_v\log N(v)=[L:\mathbf Q]\log2$, we obtain
\[
 h_F(J)-\sum_i h_F(E_i)
 =\left(\frac74-\frac32\right)\log2.
\]
All three elliptic quotients have $j=1728$, so their stable heights
are equal. This proves (JAC.P10.1).
\end{proof}

\paragraph{What this does and does not use.}
This proof computes a Hodge-lattice determinant.
It does not infer a node thickness from the unweighted intersection
graph, and it does not use the claim in Gu\`ardia's Theorem~12 sketch
that every normalized prime-power isogeny height difference has an
integer coefficient of $\log p$. The determinant calculation is
compatible with arbitrary further ramification, since all valuations
were normalized by $\nu(2)=1$.

\paragraph{Primary source locators.}
Liu's item is \emph{Example}~10.4.39, pp.~548--549 of the 2006
paperback, with the elliptic coordinate computation referred to
Example~10.1.28. The full primary book text at these pages was read
through the PDF mirror
\url{https://rexresearch1.com/GeometryLibrary/AlgebraicGeometryArithmeticCurvesLiu.pdf};
the author book page is
\url{https://www.math.u-bordeaux.fr/~qliu/Book/index.html}.
R\'emond's primary author PDF is
\url{https://www-fourier.univ-grenoble-alpes.fr/~remond/2476.pdf}.
The new content here is the smooth tail chart (JAC.P10.3), the
integral basis (JAC.P10.8)--(JAC.P10.8b), and the resulting exact
height difference. The weaker lower-bound version and the exact
basis refinement were independently checked in this run.
